\chapter{Cost-Accounted Rho}
\label{ch:carho}
% ===================================================================

Chapter~\ref{ch:costmonad} constructed a meter in the abstract: given a
continued interactive GSLT, the cost endofunctor wraps its interaction cut
so that nothing fires without a charge, and the wrapping is a monad, and
the monad has adjunctions on either side saying what installing a meter
and internalising one respectively cost.
This chapter exhibits the construction on the running example.
The result is a rho calculus in which a token stack is a term like any
other, in which authority is carried by signatures rather than by
position, and in which a computation is not run to completion and billed
afterwards but \emph{accepted} in advance, on the strength of a linear
proof that it can pay.

The change of billing model is the substantive one, and it is worth
stating before the syntax, because everything else in the chapter is
downstream of it.

\section{Two ways to bill a computation}

The familiar arrangement, and the one the earlier F1R3FLY platform
inherited, is \emph{run-to-completion}: a deployment arrives with a
declared budget, the validator executes it, the meter ticks, and if the
budget runs out the whole thing is rolled back and the deposit
forfeited.
It is simple to state and it has two costs, one economic and one
mathematical.

The economic cost is that work is done which is then thrown away.
The mathematical cost is worse.
If execution is speculative, then two deployments which interleave may
each individually be affordable and jointly not, and deciding whether
a proposed merge of two histories is admissible requires re-running the
merge.
The merge analysis burden is quadratic in the worst case and it is the
reason concurrency and metering have historically been hard to combine.

The alternative is \emph{acceptance by linear proof}.
A deployment arrives carrying its token stacks as data.
Before any reduction, the validator checks a linear proof that the stacks
suffice for every reduction the term can perform.
Acceptance is a static judgment; execution afterwards cannot fail for
want of funds, because a term which cannot pay cannot be accepted.
Merges become admissible or not on the strength of the proofs, without
re-execution.

\begin{remark}[Why this is the right shape for the rest of the book]
\label{rmk:acceptance-shape}
The learner of Chapter~\ref{ch:sci} does not first run an experiment and
then discover it could not afford it.
It decides in advance which experiments are within its means, and the
whole force of the confinement results there depends on that decision
being made \emph{before} the interaction rather than after.
Acceptance by linear proof is what makes that a fact about the calculus
rather than a modelling convenience.
\end{remark}

\section{Syntax: stacks and signatures as terms}

The calculus has four mutually defined syntactic categories: processes,
names, signed terms, and signatures.
Processes and names are as in Section~\ref{sec:rho-running}, with one
change: what a for-comprehension continues into, and what a send carries,
is a \emph{signed} term rather than a bare process.
\begin{align}
  P, Q &\;::=\; \nil
        \;\mid\; \mathtt{for}(y \leftarrow x)\,T
        \;\mid\; x!(U)
        \;\mid\; P \Par Q
        \;\mid\; {*}x \\
  x, y &\;::=\; @T \;\mid\; s
\end{align}
A signed term is a process under a signature, a parallel composition of
signed terms, or a token stack:
\begin{align}
  T, U &\;::=\; \{P\}_{s} \;\mid\; T \Par U \;\mid\; S \\
  S    &\;::=\; () \;\mid\; s : S
\end{align}
and a signature is a ground key drawn from a cryptographic backend $G$, a
hash of a process, or the composition of two signatures:
\[
  s(G) \;::=\; g \;\mid\; \#P \;\mid\; s \ast s .
\]

Three features deserve comment.

\paragraph{Token stacks are terms.}
A stack $s_1 : s_2 : \cdots : ()$ is a sequence of signature-indexed
balance layers, and it is a term of the calculus --- it can be sent,
received, quoted, and stored on a channel.
This is not a convenience.
It is what makes a purse a \emph{located} object, and located purses are
what Chapter~\ref{ch:cigslt} needed in order for authority to be a fact
about where a capability sits rather than a global ambient permission.
The empty stack is a depleted balance, and a term holding the empty stack
is a computation that has stopped.

\paragraph{There are two axes of reflection.}
Structural quoting $@T$ turns a signed term into a channel name and
preserves everything, so that ${*}x$ recovers the term including its
signature and its cost provenance.
Cryptographic quoting $\#P$ turns a process into a signature and is
one-way.
Both are reflection in the sense that they take a computational entity and
return a name; they differ in whether the entity can be got back.
The book uses the first constantly --- crossover in
Chapter~\ref{ch:com} operates on structurally quoted code --- and the
second only where authority is at issue.

\paragraph{Parallel composition is parametrically polymorphic.}
The same $\Par$ appears at the process level and at the signed-term
level.
This is not an overloading: it is the statement that composing two
metered computations is the same operation as composing two unmetered
ones, and it is what allows the cost functor of
Chapter~\ref{ch:costmonad} to be a functor at all rather than a
re-presentation of the calculus.

\section{The rules, and where the charge goes}

The pure calculus has one rule.
The metered calculus has five, and the multiplication is entirely
accounted for by the fact that a communication now has to say what
happened to the money.
Informally:

\begin{enumerate}[itemsep=3pt]
  \item \textbf{Comm.}
    A send meets a matching receive.
    The payload substitutes, and the head cell of the receiver's stack is
    consumed.
    This is the interaction cut of Chapter~\ref{ch:cigslt} with the meter
    installed on it.
  \item \textbf{Join.}
    Several sends meet a receive that requires all of them.
    The charge is not the sum of the individual charges; joins are priced
    for the synchronization, which is the only place in the calculus
    where the cost of an interaction is not local to a single cut.
  \item \textbf{Drop.}
    ${*}x$ recovers a signed term from a name, at a charge proportional to
    what is recovered.
  \item \textbf{Struct.}
    Structural congruence is free.
    Rearranging a term costs nothing, which is the formal content of the
    claim that the equations $\eqs$ are gauge.
  \item \textbf{Par.}
    Reduction propagates under composition, with the stacks of the
    components independent.
\end{enumerate}

The rule that matters most for the rest of the book is the guarded form
of comm.
A for-comprehension may carry a \texttt{where} clause,
\[
  \mathtt{for}(y \leftarrow x)\ \mathtt{where}\ \varphi \ \{ P \}_s ,
\]
in which $\varphi$ is a formula of the context-decorated modal logic of
Chapter~\ref{sec:lts-hml}.
The communication fires only if the payload satisfies $\varphi$.
This is the point at which the generated logic stops being a device for
reasoning \emph{about} programs and becomes a device programs use, and
it is the single feature the ecology chapters lean on hardest: a
learner's hypothesis is a formula, and testing a hypothesis is guarding a
communication on it.

\begin{remark}[Checkability, not adequacy]
\label{rmk:carho-checkable}
A \texttt{where} clause must be decided in bounded time on a particular
payload.
That is a different demand from the one adequacy makes, which is that the
logic as a whole distinguish exactly what the theory distinguishes.
Chapter~\ref{ch:hypercube} separates the two demands and shows that the
spatial and modal fragment suffices for the first; the discussion there
is the reason the guard is restricted to that fragment rather than to
the full logic.
\end{remark}

\section{From phlogiston to a spectrum}

In the platform's earlier design there was one kind of resource, called
phlogiston, and a single number measured it.
The cost-accounted calculus replaces the number with a vector: a stack
carries balances indexed by signature, and different rewrite rules debit
different components.

The reason is not bookkeeping fastidiousness.
A single number presumes that every kind of resource a computation
consumes is interconvertible at a fixed rate, and that presumption is
exactly what Chapter~\ref{ch:vtok} shows to be a nontrivial condition
rather than a definition.
Storage, computation, bandwidth, and authority are different
commodities.
Whether they admit a common unit is a question about the conversion
structure available, and in general the answer is no.
Decomposing phlogiston into a spectrum is what makes the question
askable.

\begin{remark}[The correction from located stacks]
\label{rmk:located-correction}
An earlier presentation allowed a term to draw on any stack in scope.
That is ambient authority in the sense of
Chapter~\ref{ch:cigslt}, and it re-admits precisely the leak the
capability discipline was introduced to close: a subterm could spend
funds it was never handed.
The correction is that a purse is located --- a stack sits at a name, and
drawing on it requires holding that name.
The consequence for the ecology chapters is not incidental.
It is what makes eating a \emph{local} act, one computation breaking
another open at a place, rather than a redistribution performed by a
scheduler standing outside the world.
\end{remark}

\section{Overcharge and refund}

One awkwardness remains, and it is worth naming because it recurs
whenever a static discipline meets data-dependent behaviour.
The linear proof required for acceptance must bound the cost of every
reduction the term can perform.
When the cost depends on data not available until run time --- the size
of a received payload, say --- the bound must be taken over the worst
case.
The term is therefore overcharged at acceptance, and the difference
refunded on completion.

The refund is not a patch.
It is the visible trace of the gap between what can be proved in advance
and what actually happens, and the gap is the same one the learner faces
in Chapter~\ref{ch:sci} when it must budget for an experiment whose
informativeness it cannot know until the experiment is done.
A framework in which acceptance were exact would be one in which
prediction were free.

% ===================================================================
